% Zonnon 
\subsection{Metrics Evaluation for Zonnon}

Zonnon is a statically-typed, object-oriented language in the Pascal/Modula-2/Oberon lineage, designed with an emphasis on abstraction unification, concurrency, and modular composition \cite{gutknecht2004zonnon}. Its object model is centered around four primary constructs—\textit{module}, \textit{object}, \textit{definition}, and \textit{implementation}—which collectively define both structural and behavioral aspects of programs. The language explicitly separates interface specification (\textit{definitions}) from implementation and supports concurrency via active objects and protocol-controlled interactions. In this section, we compute the proposed object-model metrics based strictly on the formal language specification.

\subsubsection{Type System and Subtyping Model}

\textbf{1) Nominal–Structural Typing Score $S_{nom}(L)$}

Zonnon employs a \textbf{nominal type system}. Subtyping is explicitly defined through the \texttt{implements} relation between objects and definitions: “an object implementing a definition is required to implement all of its fields and methods”. There is no notion of structural subtyping in the specification.

According to the language specification, Zonnon employs an object model where types (objects) explicitly declare conformance to interfaces (``definitions'') via the \texttt{implements} clause, and relationships such as \texttt{implements} and \texttt{refines} are explicitly stated rather than inferred structurally; hence, it satisfies nominal typing requirements with explicit subtype declarations \cite{gutknecht2004zonnon}. 

Conversely, there is no evidence of structural subtyping based solely on shape compatibility or formal structural typing rules (i.e., no typing judgments allowing implicit compatibility based on member structure). 

\begin{itemize}
    \item $I_{nom}(Zonnon) = 1$
    \item $I_{struct}(Zonnon) = 0$
\end{itemize}

\[
    S_{nom}(Zonnon) = \frac{1 + 0}{2} = 0.5
\]

\textbf{2) Variance Formalization Score $S_{var}(L)$}

Zonnon does not support parametric polymorphism (generics) or type constructors with variance annotations; instead, it relies on object types, definitions (interfaces), and explicit implementation relations without any mechanism for parameterized type abstraction. Consequently, variance cannot be expressed at the language level.
Furthermore, since variance is not part of the type system, there are no formal rules, constraints, or proofs of soundness related to variance in the specification.

\begin{itemize}
    \item $V_{exp}(Zonnon) = 0$
    \item $V_{sound}(Zonnon) = 0$
\end{itemize}

\[
S_{var}(Zonnon) = 0
\]

\textbf{3) Inheritance–Subtyping Separation Score $S_{rel}(L)$}

First, Zonnon explicitly supports subtyping via \emph{definitions} (interfaces) that are independent constructs from objects, allowing an object to implement multiple definitions without relying on inheritance hierarchies; this constitutes explicit subtyping without inheritance constructs. 

\[
    I_{dec}(Zonnon)=1
\]

Second, subtyping semantics are defined through the “implements” relation between objects and definitions, which is orthogonal to object construction and does not derive from inheritance semantics, demonstrating full independence. 

\[
    I_{ind}(Zonnon)=1
\]

Third, the specification formally distinguishes constructs such as \emph{object}, \emph{definition}, and their relations (e.g., “implements”, “refines”), thereby providing an explicit and formal separation between interface-based subtyping and structural composition.

\[
    I_{form}(Zonnon)=1
\]

Therefore

\[
    S_{rel}(Zonnon) = \frac{1 + 1 + 1}{3} = 1 
\]

\textbf{Final metric value:}

\[
    TSSI(Zonnon) = 0.3 \cdot S_{nom}(L) + 0.3 \cdot S_{var}(L) + 0.4 \cdot S_{rel}(L) = 0.3 \cdot 0.5 + 0.3 \cdot 0 + 0.4 \cdot 1 
\]

\subsubsection{Abstraction and Encapsulation Mechanisms (AE)}

\textbf{1) Visibility Granularity Index $V(L)$}

Zonnon provides explicit visibility modifiers \texttt{private} and \texttt{public}, where \texttt{private} restricts access to the local scope of the declaration and \texttt{public} allows visibility across imported program units; additionally, there exists an implicit module-level encapsulation boundary, since declarations are only accessible if explicitly exported and imported \cite{gutknecht2004zonnon}. Thus, the language supports three semantically distinct access scopes: local (private), module/exported (public), and non-exported encapsulated scope. 

\begin{itemize}
    \item $n_{levels}(Zonnon)=3$
\end{itemize}

\[
    V(Zonnon)=\frac{3}{n_{max}}
\]

\textbf{2) Module Isolation Strength $M(L)$}

The language defines four primary program units: \texttt{module}, \texttt{object}, \texttt{definition}, and \texttt{implementation}, which together form the basis of program composition and encapsulation. 
Among these, \texttt{module}, \texttt{definition}, and \texttt{object} enforce strict export control via explicit use of visibility modifiers (e.g., \texttt{public}) and require explicit \texttt{import} declarations for external access, thereby acting as boundaries with enforced interface exposure. 
In contrast, \texttt{implementation} serves as a reusable code container aggregated into other units and does not independently enforce a strict export boundary.

\[
M(Zonnon)=\frac{3}{4} = 0.75
\]

\textbf{3) Interface Abstraction Capacity $I(L)$}

The language defines \texttt{definition} as a dedicated abstraction construct representing an interface (a set of method signatures and fields), clearly decoupled from implementation, while \texttt{implementation} provides reusable method bodies, and \texttt{object} serves as the concrete type that implements one or more definitions via the \texttt{implements} clause \cite{gutknecht2004zonnon}. 
Additionally, \texttt{protocol} types describe interaction behavior but are primarily oriented toward concurrency rather than general interface abstraction, and thus are not counted as core abstraction constructs.

\begin{itemize}
    \item $n_{abstraction\_constructs}(Zonnon)=1$
    \item $n_{object}(Zonnon)=2$
\end{itemize}

\[
I(Zonnon)=\frac{1}{1}=1.0
\]

\textbf{4) Encapsulation Enforcement Coefficient $E(L)$}

The specification enforces encapsulation through visibility modifiers (\texttt{private}, \texttt{public}) and requires all access to object state to occur via declared fields or methods, with no indication of unchecked field access or unsafe casts that bypass the type system. 
Although Zonnon includes a reflection facility, it is defined in a controlled manner based on structured metadata (e.g., XML schema) and does not provide arbitrary runtime mutation or violation of access restrictions. Furthermore, type conversions and interface checks (e.g., \texttt{is}, \texttt{implements}) are explicitly constrained and type-safe.

\[
    E(Zonnon) = 1
\]

\textbf{Final metric value:}

\[
    AE(Zonnon) = 0.2 \cdot V(L) + 0.25 \cdot M(L) + 0.25 \cdot I(L) + 0.3 \cdot E(L)
\]

\subsubsection{Inheritance Semantics and Hierarchy Management}

\textbf{1) Hierarchy Simplicity Factor $H_s(L)$}

Zonnon does not support classical multiple inheritance; reuse is achieved via implementations and definitions.
According to the specification, Zonnon does not support class-based inheritance between \texttt{object} types; objects cannot extend other objects and thus there is no notion of multiple superclass inheritance. Instead, reuse is achieved via \texttt{implementation} aggregation and definition conformance via \texttt{implements}, which are not counted as superclass inheritance.

\begin{itemize}
    \item $M(Zonnon)=0 \Rightarrow H_s(Zonnon)=1$
\end{itemize}

\textbf{2) Linearization Determinism Score $L_d(L)$}

The language does not support class inheritance hierarchies and therefore does not require a method resolution order (MRO) for multiple inheritance. Instead, behavior composition is achieved through explicit \texttt{implements} relations with \texttt{definition}s and optional aggregation from \texttt{implementation} units, where method bindings are statically defined and unambiguous at compile time \cite{gutknecht2004zonnon}. Since there is no inheritance-based method overriding chain requiring linearization, and all method resolutions are explicitly specified and deterministic by construction, the language effectively satisfies the strongest determinism condition.

\[
    L_d(Zonnon)=1
\]


\textbf{3) Constraint Protection Coefficient $C_p(L)$}

$p_1=0$: Zonnon does not provide explicit \texttt{final} or non-inheritable class/method modifiers. 

$p_2=0$: Zonnon does not support sealed or closed class hierarchies with explicitly enumerated subclasses.

$p_3=0$: There is no mandatory \texttt{override} annotation enforced by the compiler for method redefinition 

$p_4=1$: Zonnon follows a model where methods are not implicitly virtual unless defined within extensible class  contexts, effectively aligning with non-virtual-by-default semantics 

$p_5=1$: Zonnon enforces a clear distinction between abstract definitions (via definition modules and abstract  classes) and implementations, supporting separation constraints 

$p_6=1$: Zonnon includes compile-time inheritance restrictions such as controlled visibility and module-based encapsulation that limit illegal extensions.

\[
    C_p(Zonnon) = \frac{3}{6} = 0.5
\]


\textbf{4) Fragile Base Mitigation Factor $F_b(L)$}

$s_1=0$: Zonnon does not require mandatory \texttt{override} annotations.

$s_2=1$: Zonnon enforces strict visibility control through module-based encapsulation and access modifiers, limiting unintended exposure of overridable members.

$s_3=0$: Zonnon does not mandate explicit superclass invocation (e.g., enforced \texttt{super} calls).

$s_4=1$: Zonnon relies on a statically typed system with type-safe method redefinitions, implying basic variance safety. 

$s_5=1$: The methods are effectively non-overridable by default unless defined within extensible class contexts, corresponding to opt-in polymorphism.

$s_6=1$: The compiler enforces consistency of method signatures during overriding, though without formal contract checking.

\[
    F_b(Zonnon) = \frac{4}{6} \approx 0.67
\]

\textbf{Final metric value:}

\[
    ISRI(Zonnon) = 0.2 \cdot H_s(L) + 0.25 \cdot L_d(L) + 0.3 \cdot C_p(L) + 0.25 \cdot F_b(L)
\]

\subsubsection{Composition and Alternatives to Inheritance}

\textbf{1) Class-Based Delegation Score $D_{cls}(L)$}

The language does not provide dedicated syntactic delegation keywords (such as \texttt{delegate} in some languages); however, it enables composition and reuse through \texttt{implementation} aggregation and object composition, which can be used to emulate delegation patterns manually. 
Zonnon does not include mechanisms for automatic forwarding of method calls (i.e., no language-level support for implicit delegation without boilerplate code).

\begin{itemize}
    \item Belegation is supported only indirectly $\Rightarrow K_{del}=0.5$
    \item All forwarding must be explicitly implemented $\Rightarrow S_{fwd}=0$
\end{itemize}

\[
D_{cls}(Zonnon)=0.25
\]

\textbf{2) Trait Expressiveness Score $T_{exp}(L)$}

Zonnon does not provide traits or mixins as first-class language constructs; instead, it supports composition via \texttt{implementation} aggregation and interface conformance via \texttt{definition}. In cases where multiple implementations introduce conflicting method names, the specification requires explicit resolution at the object level, and there is no support for aliasing or exclusion mechanisms; thus, conflicts result in a compilation requirement for disambiguation. 
Additionally, \texttt{implementation} units do not declare behavioral requirements on consuming objects (unlike traits with required methods); they are purely containers of reusable code without constraint specification.

\[
    T_{exp}(Zonnon)=0
\]

\textbf{3) Interface Composition Capacity $I_{comp}(L)$}

In Zonnon, \texttt{object} types can implement multiple \texttt{definition}s (interfaces) via the \texttt{implements} clause without any specified upper bound, implying unrestricted interface composition \cite{gutknecht2004zonnon}. 
However, \texttt{definition} constructs serve purely as abstract interfaces containing field declarations and method signatures; they do not allow method implementations, as behavior is provided separately via \texttt{implementation} or directly within \texttt{object}s.

\begin{itemize}
    \item $N_{impl}(Zonnon)=1$
    \item $D_{def}(Zonnon)=0$
\end{itemize}

\[
    I_{comp}(Zonnon)=0.5
\]

\textbf{Final metric value:}

\[
    CRFI(Zonnon) = 0.3 \cdot D_{cls}(L) + 0.35 \cdot T_{exp}(L) + 0.35 \cdot I_{comp}(L)
\]

\subsubsection{Object Representation and Creation Model}

\textbf{1) Identity Semantics Score $S_{id}(L)$}

According to the specification, the primary user-definable types are \texttt{object} and \texttt{record}. Objects in Zonnon can be declared with modifiers \texttt{ref} (reference semantics) or \texttt{value} (value semantics), with \texttt{value} being the default, while \texttt{record} is explicitly defined as a value object type \cite{gutknecht2004zonnon}. Since reference semantics are not enforced uniformly across all user-defined types—only optionally for \texttt{object} types—there is no type category that strictly enforces reference semantics by definition.

Zonnon distinguishes value and reference objects:

\begin{itemize}
    \item $T_{user}=2$ (value, reference)
    \item $T_{ref}=1$
\end{itemize}

\[
    S_{id}(Zonnon)=0.5
\]

\textbf{2) Instantiation Paradigm Diversity $S_{cr}(L)$}

The specification defines instantiation primarily via the \texttt{new} operator, which is used to create instances of \texttt{object} types and dynamically allocated arrays. Additionally, Zonnon supports the instantiation of concurrent entities through activity creation (i.e., spawning \texttt{activity} instances with their own execution thread), representing a distinct actor-like paradigm \cite{gutknecht2004zonnon}. However, the language does not support prototype-based cloning or object literal construction for user-defined objects. 

\[
S_{cr}(Zonnon) = \frac{2}{P_{max}}
\]

\textbf{3) Meta-level Extensibility Score $S_{meta}(L)$}

The specification defines a reflection mechanism that enables querying program structure and metadata (e.g., via XML schema representations), allowing inspection of types, declarations, and program composition at runtime \cite{gutknecht2004zonnon}. However, there is no support for modifying object structures dynamically, nor for altering the instantiation process or introducing metaobject protocols. Thus, Zonnon provides introspection without structural or behavioral intercession.

\[
    S_{meta}(Zonnon) = \frac{1}{I_{max}}
\]

\textbf{Final metric value:}

\[
    ORCI(Zonnon) = 0.35 \cdot S_{id}(L) + 0.3 \cdot S_{cr}(L) + 0.35 \cdot S_{meta}(L)
\]

\subsubsection{Polymorphism and Dispatch Mechanisms}

\textbf{1) Dispatch Scope Mechanism $M_{scope}(L)$}

Zonnon supports single dispatch, where method invocation is determined by the dynamic type of the receiver object (i.e., receiver-based method binding) \cite{gutknecht2004zonnon}. 

$I_{single}=1$

However, the language does not support multiple dispatch, as method selection is not based on the runtime types of multiple arguments.

$I_{multi}=0$

There is no support for predicate-based or condition-based dispatch mechanisms (such as guard-based method selection); control flow must be expressed explicitly within method bodies.

$I_{pred}=0$

\[
    M_{scope}(Zonnon)=\frac{1}{5}=0.2
\]

\textbf{2) Dispatch Safety Guarantee $S_{disp}(L)$}

According to the specification, method calls are statically type-checked against \texttt{definition}s (interfaces), and \texttt{object}s must fully implement all declared methods before instantiation; thus, method availability is guaranteed at compile time \cite{gutknecht2004zonnon}. Furthermore, there is no support for unsafe downcasting or dynamic method lookup that could lead to “method not found” errors during dispatch. Type tests (e.g., \texttt{is}) are explicitly checked and must be handled by the programmer, preventing invalid dispatch paths.

First, Zonnon is a statically typed language with compile-time type checking of method calls and dynamic dispatch resolution constrained by class hierarchies. Incorrect method calls are rejected at compile time.

$I_{static}(\text{Zonnon})=1$

Second, the language does not enforce exhaustive dispatch definitions: there is no requirement that all possible variants of a type hierarchy must be covered (e.g., no pattern matching exhaustiveness or sealed hierarchy guarantees), and method overriding is optional rather than complete by construction.

$I_{exhaust}(\text{Zonnon})=0$


\[
    S_{disp}(Zonnon)=\frac{1 + 0}{2} = 0.5
\]

\textbf{3) Virtualization Optimization Potential $O_{virt}(L)$}

Zonnon does not support classical method overriding or virtual dispatch; all methods are effectively non-virtual due to the absence of class inheritance and overriding semantics, which corresponds to implicit support for final methods \cite{gutknecht2004zonnon}.

$K_{final}=1$  

The language also does not define class hierarchies and therefore does not include constructs such as \texttt{sealed} for restricting subclassing.

$K_{sealed}=0$

Zonnon supports procedures and functions that are statically bound (i.e., not subject to dynamic dispatch), which serves as an inherent static dispatch mechanism.

$K_{static}=1$

\[
    O_{virt}(Zonnon)=\frac{2}{3} \approx 0.67 
\]

\textbf{Final metric value:}

\[
    PDEI(Zonnon) = 0.35 \cdot M_{scope}(L) + 0.35 \cdot S_{disp}(L) + 0.3 \cdot O_{virt}(L)
\]

\subsubsection{State Model, Mutability, and Concurrency Guarantees}

\textbf{1) Default Mutability Score $M_{def}(L)$}

The \texttt{object} types are mutable by default, allowing modification of their fields unless explicitly restricted by qualifiers such as \texttt{immutable} or by using value semantics (e.g., \texttt{value} objects or \texttt{record}s) \cite{gutknecht2004zonnon}. 

\[
    M_{def}(Zonnon) = 0
\]

\textbf{2) Static Reference Constraint Density $S_{dens}(L)$}

The language provides reference semantics primarily through \texttt{ref} objects and implicitly via object references, which constitute reference-related constructs alongside qualifiers such as \texttt{immutable} that affect reference behavior \cite{gutknecht2004zonnon}. 

$N_{ref}(Zonnon)=2$ (reference objects via \texttt{ref}, and reference qualifiers such as \texttt{immutable}). 

Among these, only \texttt{immutable} imposes a static constraint by restricting mutation across scope boundaries at compile time, while \texttt{ref} itself introduces reference semantics without additional safety guarantees.

First, Zonnon includes syntactic constructs that introduce or manipulate references: object variables declared with \texttt{var}, which denote references to objects, class types, since all class instances are accessed via references, and procedure parameters, which may implicitly pass references depending on their type semantics.

$N_{ref}(\text{Zonnon}) = 3$. 

Second, Zonnon includes only constructs that impose static constraints on reference usage: the \texttt{immutable} modifier, which enforces compile-time restrictions on mutation through references, and access modifiers (e.g., \texttt{protected}, \texttt{public}) that constrain reference accessibility at compile time.

$N_{constr}(\text{Zonnon}) = 2$

\[
    S_{dens}(Zonnon) = \frac{2}{3} = \approx 0.67
\]

\textbf{3) Concurrency Primitive Safety $P_{conc}(L)$}

The language introduces concurrency through \texttt{activity} constructs, which encapsulate independent threads of execution and communicate via well-defined interfaces and synchronization mechanisms, resembling an actor-like model \cite{gutknecht2004zonnon}. However, the specification does not provide formal guarantees of data race freedom enforced by the type system; shared state may still be accessed, and correctness relies on disciplined use of synchronization constructs.

$S_{race}(Zonnon)=0$

On the other hand, the active object model represents a higher-level abstraction compared to raw threads and locks, as it structures concurrency around object interaction and implicit synchronization, aligning partially with actor-like paradigms but without strict isolation guarantees; thus, it corresponds to structured/high-level primitives.

$S_{model}(Zonnon)=0.5$

\[
    P_{conc}(Zonnon) = \frac{0 + 0.5}{2} = 0.5
\]

\textbf{Final metric value:}

\[
    SCSI(Zonnon) = 0.35 \cdot M_{def}(L) + 0.25 \cdot S_{dens}(L) + 0.4 \cdot P_{conc}(L) = 0.35 \cdot 0 + 0.25 \cdot 0.67 + 0.4 \cdot 0.5
\]

\subsubsection{Support for Formal Specification and Verifiability}

\textbf{1) Native Contract Constructs Score $C_{nat}(L)$}

$n_{pre}=1$: Zonnon includes explicit \texttt{require} clauses in procedure declarations, which define preconditions and are both syntactically integrated and semantically enforced.

$n_{post}=1$: Zonnonalso provides \texttt{ensure} clauses for postconditions, similarly defined at the language level.

$n_{inv}=0$: Zonnon does not define a dedicated invariant construct (e.g., class invariants) as a first-class syntactic feature. Invariants must be expressed indirectly through assertions or programmer discipline.


\[
    C_{nat}(Zonnon) = \frac{2}{3} = \approx 0.67
\]

\textbf{2) Verification Enforcement Level $V_{enf}(L)$}

Zonnon does not provide a static verification system capable of proving correctness properties at compile time; there is no integrated formal verification or proof-based type system. However, the language defines preconditions and postconditions with a well-specified runtime semantics, meaning these conditions can be checked during program execution (e.g., triggering exceptions if violated). Therefore, enforcement exists but is limited to runtime checking rather than compile-time guarantees.

\[
    V_{enf}(Zonnon) = 0.5
\]

\textbf{Final metric value:}

\[
    FSVI(Zonnon) = 0.45 \cdot C_{nat}(L) + 0.55 \cdot V_{enf}(L) = 0.45 \cdot 0.67 + 0.55 \cdot 0.5
\]

